\section{Covering a tree}\label{sec:trees}

We now cover an entire chain of extensions. The input is a finite tree
whose edge labels are disjoint along each path and whose child laws are
spread. The aim is to put the labels of one full path inside a product
sample, with a cost logarithmic in the largest extension size.

A \emph{probability tree on $X$} is a finite rooted tree with all
leaves at the same depth, together with a probability distribution
$\nu_v$ on the children $\ch(v)$ of each non-leaf node $v$ and a label
$A_v(b)\subseteq X$ on the arc from $v$ to each child $b$.
We delete zero-weight children. Different nodes or arcs may carry
identical labels; nodes are never identified merely because their
labels agree. A \emph{path} always means a full root-to-leaf path,
unless otherwise specified. Write
\[
 \UU(P)=\bigcup_{(v,b)\in P}A_v(b)
\]
for its union of labels.

The tree is \emph{disjoint} if the labels along each path are pairwise
disjoint. A level has \emph{capacity} $m$ if all its labels have size
at most $m$. A level is \emph{$a$-spread} if the child distribution
at every node immediately above that level is $a$-spread.
Thus all requirements concern individual child distributions;
no spread assumption is made on the distribution of a full path.

\Needspace{13\baselineskip}
\begin{theorem}[Tree covering]\label{thm:tree}
Let $\TT$ be a disjoint probability tree on a finite set $X$,
with $k\ge1$ levels. Suppose that every level is $\sigma$-spread,
where $\sigma>0$, and that its positive integer capacities satisfy
\[
 \ell_{i+1}\ge16\ell_i\qquad(1\le i<k).
\]
Set
\[
 r=\min\{1,\;16\ee^{50}\sigma(1+\log_2\ell_k)\}.
\]
Then
\[
 \PP_{W\sim\mu_r}
 \bigl(\UU(P)\subseteq W\text{ for some path }P\text{ of }\TT\bigr)
 \ge\tfrac23.
\]
\end{theorem}

Section~\ref{sec:graphs} constructs such a tree for $H$ with
$\sigma=128\pE(n,H)$ and $\ell_k\le e(H)$. We prove the tree theorem
by repeatedly pruning and reweighting the tree while replacing its
labels by smaller subsets. Empty labels and unequal label sizes
must therefore be allowed.

\subsection{One simultaneous reduction}

\begin{lemma}[Reduction]\label{lem:reduction}
Let $\TT$ be a disjoint probability tree with $k\ge1$ levels and positive integer
capacities $m_1,\ldots,m_k$ satisfying $m_{j+1}\ge16m_j$.
Suppose level $j$ is $\alpha_j$-spread, where
\[
 0<\alpha_j\le a,\qquad \rho=\ee^{50}a<1.
\]
There is a deterministic procedure which, given $W\sim\mu_\rho$,
succeeds with probability at least $3/4$ and on success returns a
disjoint probability tree $\TT'$ with the same levels, such that:
\begin{enumerate}[label=(\roman*),leftmargin=2.2em,itemsep=3pt]
 \item level $j$ has capacity $\lfloor m_j/2\rfloor$ and is
 $\alpha_j/(1-\ee^{-m_j})$-spread;
 \item for every path $Q$ of $\TT'$ there is a path $P$ of $\TT$ with
 \begin{equation}\label{eq:recovery}
   \UU(P)\subseteq W\cup\UU(Q).
 \end{equation}
\end{enumerate}
\end{lemma}

\begin{proof}
The output will be a subtree with labels $A_v(b)\setminus W$ and
child laws obtained from the target marginals of the local transition.
We first identify which children can be retained while discarding
at most $\ee^{-m_j}$ of the transition mass at each retained node
with outgoing level $j$.
Keep the input tree fixed throughout this analysis. At a node whose
outgoing level is $j$, use the transition of
Section~\ref{sec:resampling} with its child distribution and
capacity $m_j$. All nodes use the same set $W$. At fixed $W$, write
$d_v(W)$ for the probability of local failure at $v$, and
$t_{vb}(W)$ for the probability of local non-failure with target $b$.
Thus
\begin{equation}\label{eq:local-masses}
 d_v(W)+\sum_{b\in\ch(v)}t_{vb}(W)=1.
\end{equation}
All these probabilities refer to the child laws of the fixed input
tree, before any pruning or conditioning.

\medskip\noindent
\textit{Pruning.}
For the remainder of the construction, fix $W$ and put
$\epsilon_j=\ee^{-m_j}$. Declare all leaves good. Working upward,
declare a node $v$ with outgoing level $s$ good if
\begin{equation}\label{eq:good-node}
 \sum_{\substack{b\in\ch(v)\\b\ \mathrm{good}}}t_{vb}(W)
 \ge1-\epsilon_s.
\end{equation}
The attempt succeeds exactly when the root is good.

To bound the probability that the root is bad, define a nonnegative
potential on every node of the input tree. Set $D_v(W)=0$ at leaves
and, at a node with outgoing level $s$, define recursively
\begin{equation}\label{eq:pruning-recurrence}
 D_v(W)=\frac{d_v(W)+\sum_{b\in\ch(v)}t_{vb}(W)D_b(W)}{\epsilon_s}.
\end{equation}
Every bad node satisfies $D_v(W)>1$. Indeed, assuming this for its
bad children, the numerator in \eqref{eq:pruning-recurrence} is at
least
\[
 d_v(W)+\sum_{b\ \mathrm{bad}}t_{vb}(W)>\epsilon_s,
\]
where the last inequality follows from \eqref{eq:local-masses} and
the failure of \eqref{eq:good-node}. This proves the assertion by
backward induction.

\medskip\noindent
\textit{Expected potential.}
We now average over $W\sim\mu_\rho$. Fix a child $b$ of a node $v$
with outgoing level $s$, and abbreviate its incoming label to $A_b$.
Every label in the entire subtree below $b$ is disjoint from $A_b$:
each such label lies on a path containing the arc into $b$.
The transition at any node in that subtree only inspects $W$ on the union
of that node's child labels. By its recursive definition, $D_b(W)$
therefore depends only on $W$ outside $A_b$.

This statement concerns every node of the fixed input tree, before
conditioning on that node being good or being retained. Figure~\ref{fig:outside}
shows the disjointness that makes this possible.
\begin{figure}[!ht]
\centering
\begin{tikzpicture}[x=1cm,y=1cm,>=Stealth,
  vertex/.style={circle,draw,fill=white,inner sep=2pt},
  every node/.style={font=\small}]
\node[vertex] (v) at (0,0) {$v$};
\node[vertex] (b) at (-1.8,-1.0) {$b$};
\node[vertex] (c) at (1.8,-1.0) {$c$};
\node[vertex] (b1) at (-2.6,-2.0) {$b_1$};
\node[vertex] (b2) at (-1.0,-2.0) {$b_2$};
\draw[->,thick,blue!65!black] (v)--node[above left] {$A_v(b)$} (b);
\draw[->] (v)--node[above right] {$A_v(c)$} (c);
\draw[->] (b)--(b1);
\draw[->] (b)--(b2);
\draw[densely dotted] (b1)--(-2.9,-2.6);
\draw[densely dotted] (b1)--(-2.3,-2.6);
\draw[densely dotted] (b2)--(-1.3,-2.6);
\draw[densely dotted] (b2)--(-0.7,-2.6);
\node[align=left,anchor=west] at (0.2,-2.0)
  {Every label below $b$ avoids $A_v(b)$.\\
   Hence $D_b(W)$ uses only $W\setminus A_v(b)$.};
\node[anchor=south] at (0,0.4) {The same set $W$ is used at every node.};
\end{tikzpicture}
\caption{Disjointness along paths controls the change of measure. Every
arc below $b$ lies on a path through its incoming arc, so the potential
$D_b$ is a function of coordinates outside the incoming label. The diagram
is schematic; labels on different branches may overlap.}
\label{fig:outside}
\end{figure}


In the local transition at $v$,
\[
 t_{vb}(W)\le
 \PP\bigl(\mathrm{target}=b,\ Z(Y)\ge\ee^{-10m_s}\mid W\bigr).
\]
We may consequently apply \eqref{eq:distortion} with $f(W)=D_b(W)$.
Together with \eqref{eq:local-failure} and
$\epsilon_s^{-1}=\ee^{m_s}$, the recurrence gives
\begin{equation}\label{eq:expected-potential}
 \EE D_v(W)
 \le \ee^{-8m_s}
    +\ee^{12m_s}\sum_{b\in\ch(v)}\nu_v(b)\,\EE D_b(W).
\end{equation}
This step controls the dependence caused by reusing $W$.
Since the child weights sum to one, backward induction yields
\[
 \EE D_v(W)\le
 \sum_{j=s}^k
 \exp\left(12\sum_{h=s}^{j-1}m_h-8m_j\right).
\]
In particular, if $o$ is the root, then
\begin{align*}
 \PP(o\text{ is bad})
 &\le\EE D_o(W)\\
 &\le\sum_{j=1}^k
   \exp\left(12\sum_{h<j}m_h-8m_j\right)
 <\tfrac14.
\end{align*}
For the last inequality, geometric growth gives
$\sum_{h<j}m_h\le m_j/15$, and the $m_j$ are distinct positive
integers. Thus the sum is at most
$\sum_{m\ge1}\ee^{-7m}<1/4$.

\medskip\noindent
\textit{The reduced tree.}
On success, start from the good root. At any good node $v$ with
outgoing level $j$, put
\[
 \gamma_v(W)=\sum_{b\ \mathrm{good}}t_{vb}(W)\ge1-\epsilon_j.
\]
Retain each good child $b$ with $t_{vb}(W)>0$, give it weight
$t_{vb}(W)/\gamma_v(W)$, replace its incoming label by
\[
 T_b=A_v(b)\setminus W,
\]
and continue recursively into $b$. The result is a pruned and
reweighted subtree of the input tree. It has the same depth because
$\gamma_v>0$ at every retained non-leaf node. Its labels are subsets
of the old labels, so it is disjoint. Since
$A_v(b)\subseteq W\cup T_b$ on every retained arc, it also satisfies
\eqref{eq:recovery}. A retained child has $t_{vb}>0$, which implies
$|T_b|\le\lfloor m_j/2\rfloor$.

This new child law is the target marginal of the joint transition
conditioned on local non-failure and a good target. Denote this
conditioning event by $E_v$. For the spread bound, use the source:
\eqref{eq:two-inclusions} gives $T\subseteq A_v(c)$ for every
transition, and the unconditioned source has law $\nu_v$ even at fixed
$W$. Thus, for every nonempty $J\subseteq X$,
\begin{align*}
 \PP(J\subseteq T\mid E_v,W)
 &\le\frac{\PP_{c\sim\nu_v}(J\subseteq A_v(c))}{\gamma_v(W)}\\
 &\le\frac{\alpha_j^{|J|}}{1-\epsilon_j}
 \le\left(\frac{\alpha_j}{1-\epsilon_j}\right)^{|J|}.
\end{align*}
This proves all the assertions. The quantities $t_{vb}(W)$ define
the entire construction deterministically; sources and targets are
auxiliary random variables used to analyze its weights.
\end{proof}

\subsection{Iteration}

\begin{proof}[Proof of Theorem~\ref{thm:tree}]
If $r=1$, the conclusion is immediate. Suppose $r<1$, and set
\[
 a=4\sigma,\qquad \rho=\ee^{50}a,
 \qquad s=\lfloor\log_2\ell_k\rfloor+1.
\]
In particular $0<\rho<1$. Start with $\TT$ and make successive
attempts using fresh independent sets of law $\mu_\rho$. A failed attempt leaves
the current tree unchanged. After a successful attempt, retain the
tree provided by Lemma~\ref{lem:reduction}.
Conditionally on the complete past, the current tree is fixed and the
next draw still has law $\mu_\rho$.

After $t$ successes, the capacity at original level $i$ is
\begin{equation}\label{eq:capacities}
 m_i(t)=\left\lfloor\frac{\ell_i}{2^t}\right\rfloor,
\end{equation}
using $\lfloor\lfloor x\rfloor/2\rfloor=\lfloor x/2\rfloor$.
Levels of zero capacity form an initial segment. After each success,
choose any one path through that segment and keep its ending
subtree, deleting the segment. All committed labels are empty, so
each path of the retained subtree extends to a full path by adjoining
the chosen empty prefix. This preserves the recovery guarantee and every
bound at a retained node. If no levels remain, only the root is left.
Fix a tie-breaking rule for these choices. The remaining positive capacities still
satisfy $m_{i+1}(t)\ge16m_i(t)$, since
$\lfloor16x\rfloor\ge16\lfloor x\rfloor$.

At a remaining original level $i$, a valid spread parameter after
$t$ successes is
\begin{equation}\label{eq:spread-iteration}
 \alpha_i(t)=\sigma
 \prod_{u=0}^{t-1}
 \left(1-\ee^{-\lfloor\ell_i/2^u\rfloor}\right)^{-1}.
\end{equation}
The positive integers $\lfloor\ell_i/2^u\rfloor$ in this product are
distinct. The infinite product over all positive integers is less
than $4$, because
\begin{align*}
 \sum_{m=1}^{\infty}-\log(1-\ee^{-m})
 &\le \frac{1}{1-\ee^{-1}}\sum_{m=1}^{\infty}\ee^{-m}\\
 &=\frac{\ee^{-1}}{(1-\ee^{-1})^2}<\log 4.
\end{align*}
Consequently $\alpha_i(t)\le4\sigma=a$. The initial tree satisfies
these bounds, and Lemma~\ref{lem:reduction} proves them inductively
after every success. Thus every attempt before completion has
conditional success probability at least $3/4$.

After $s$ successes, \eqref{eq:capacities} is zero at every level.
Composing \eqref{eq:recovery} through the successful reductions shows
that the union of their draws contains the labels of a path in the
original tree. The expected waiting time for each next success is
at most $4/3$: conditionally on the past, its tail is bounded by
that of a geometric random variable with success probability $3/4$.
Hence the expected number of attempts needed is at most $4s/3$.
By Markov's inequality, completion occurs within $4s$ attempts with
probability at least $2/3$.

Predraw $4s$ independent sets of law $\mu_\rho$, including any that
the procedure ultimately does not use, and take their union $W_* $.
It has law $\mu_{r_*}$, where
\[
 r_*=1-(1-\rho)^{4s}
 \le4s\rho
 =16\ee^{50}\sigma s
 \le16\ee^{50}\sigma(1+\log_2\ell_k)=r.
\]
With probability at least $2/3$, $W_*$ contains a path union.
The same bound holds for $\mu_r$ by monotonicity.
\end{proof}

